\documentclass[11pt]{article}

\usepackage{amsmath}
\usepackage{amsfonts}
\usepackage{graphicx}
\usepackage{hyperref}
\usepackage{physics}

\title{Mathematical Foundations of the MedVision AI Project}
\author{Yann Sofian Smatti}
\date{}

\begin{document}

\maketitle

\tableofcontents

\newpage

\section{Introduction}

The MedVision AI project aims to build a machine learning system capable of detecting medical conditions in radiographic images.

This document explains the mathematical and algorithmic foundations of the project at a level suitable for advanced undergraduate students (L2–L3 / CPGE).

The project involves the following concepts:

\begin{itemize}
\item image representation
\item convolution
\item convolutional neural networks
\item optimization algorithms
\item transfer learning
\item evaluation metrics
\end{itemize}

\section{Mathematical Representation of Images}

A digital image can be modeled mathematically as a function.

\[
I : \Omega \rightarrow \mathbb{R}
\]

where

\begin{itemize}
\item $\Omega \subset \mathbb{Z}^2$ represents pixel coordinates
\item $I(x,y)$ represents pixel intensity
\end{itemize}

For a grayscale image:

\[
I \in \mathbb{R}^{H \times W}
\]

where

\begin{itemize}
\item $H$ is the image height
\item $W$ is the image width
\end{itemize}

Each pixel value is typically in the interval

\[
I(x,y) \in [0,255]
\]

For RGB images:

\[
I \in \mathbb{R}^{H \times W \times 3}
\]

where the third dimension corresponds to color channels.

\section{Image Preprocessing}

Neural networks require normalized inputs.

Normalization transforms the pixel values:

\[
x' = \frac{x}{255}
\]

so that

\[
x' \in [0,1]
\]

This improves numerical stability during optimization.

Another important preprocessing step is resizing.

In the MedVision project images are resized to:

\[
224 \times 224
\]

because this corresponds to the input size expected by most pretrained CNN architectures.

\section{Discrete Convolution}

Convolution is the fundamental operation of convolutional neural networks.

Let

\[
I : \mathbb{Z}^2 \rightarrow \mathbb{R}
\]

be an image and

\[
K : \mathbb{Z}^2 \rightarrow \mathbb{R}
\]

be a convolution kernel.

The discrete convolution is defined as:

\[
(I * K)(x,y)
=
\sum_{i=-k}^{k}
\sum_{j=-k}^{k}
I(x-i,y-j)K(i,j)
\]

This operation produces a new image called a \textbf{feature map}.

\subsection{Interpretation}

Convolution allows the extraction of local patterns such as

\begin{itemize}
\item edges
\item textures
\item shapes
\end{itemize}

Example kernel detecting vertical edges:

\[
K =
\begin{bmatrix}
1 & 0 & -1 \\
1 & 0 & -1 \\
1 & 0 & -1
\end{bmatrix}
\]

\section{Convolutional Neural Networks}

A convolutional neural network is composed of several layers.

Typical architecture:

\begin{itemize}
\item convolution layer
\item activation function
\item pooling layer
\item fully connected layer
\end{itemize}

A convolutional layer computes

\[
h = f(W * x + b)
\]

where

\begin{itemize}
\item $W$ are learnable filters
\item $b$ is a bias term
\item $f$ is a nonlinear activation function
\end{itemize}

\section{Activation Functions}

The most common activation function is the Rectified Linear Unit.

\[
ReLU(x) = \max(0,x)
\]

Advantages:

\begin{itemize}
\item prevents vanishing gradients
\item accelerates training
\end{itemize}

\section{Pooling}

Pooling reduces spatial dimensions.

Max pooling is defined as

\[
y = \max(x_1,x_2,\dots,x_n)
\]

Example:

\[
\begin{bmatrix}
1 & 3 \\
2 & 0
\end{bmatrix}
\rightarrow 3
\]

Pooling reduces computation and improves translation invariance.

\section{Classification Layer}

After feature extraction, the network performs classification.

The output layer uses the sigmoid function:

\[
\sigma(x) = \frac{1}{1 + e^{-x}}
\]

which maps real numbers into probabilities.

\section{Loss Function}

Binary classification problems use binary cross entropy.

\[
L = -\frac{1}{N}
\sum_{i=1}^{N}
\left[
y_i \log(p_i)
+
(1-y_i)\log(1-p_i)
\right]
\]

where

\begin{itemize}
\item $y_i$ is the true label
\item $p_i$ is the predicted probability
\end{itemize}

\section{Optimization}

Parameters are updated using gradient descent.

\[
\theta_{t+1} = \theta_t - \eta \nabla L(\theta)
\]

where

\begin{itemize}
\item $\theta$ represents the model parameters
\item $\eta$ is the learning rate
\end{itemize}

The gradient is computed using backpropagation.

\section{Backpropagation}

Backpropagation relies on the chain rule.

If

\[
f(x) = f_3(f_2(f_1(x)))
\]

then

\[
\frac{df}{dx}
=
\frac{df_3}{df_2}
\cdot
\frac{df_2}{df_1}
\cdot
\frac{df_1}{dx}
\]

This allows efficient gradient computation in deep networks.

\section{Transfer Learning}

Training deep CNNs from scratch requires millions of images.

Transfer learning reuses models trained on large datasets such as ImageNet.

The model can be written as

\[
f(x) = g(h(x))
\]

where

\begin{itemize}
\item $h(x)$ is the pretrained feature extractor
\item $g(x)$ is the new classifier
\end{itemize}

In the MedVision project the feature extractor is frozen and only the final layers are trained.

\section{Data Augmentation}

Medical datasets are often small.

Data augmentation generates new samples using transformations:

\begin{itemize}
\item rotation
\item flipping
\item zoom
\item brightness modification
\end{itemize}

Mathematically:

\[
x' = T(x)
\]

where $T$ is a transformation operator.

\section{Evaluation Metrics}

In medical diagnostics accuracy alone is insufficient.

We use several metrics.

\subsection{Accuracy}

\[
Accuracy =
\frac{TP + TN}{TP + TN + FP + FN}
\]

\subsection{Precision}

\[
Precision =
\frac{TP}{TP + FP}
\]

\subsection{Recall}

\[
Recall =
\frac{TP}{TP + FN}
\]

\subsection{F1 Score}

\[
F1 =
2
\frac{Precision \cdot Recall}
{Precision + Recall}
\]

\section{Explainable AI}

Medical models must be interpretable.

Grad-CAM produces a heatmap highlighting regions used for prediction.

The Grad-CAM map is defined by

\[
L^c =
ReLU
\left(
\sum_k
\alpha_k^c A^k
\right)
\]

where

\begin{itemize}
\item $A^k$ are feature maps
\item $\alpha_k^c$ are weights computed using gradients
\end{itemize}

\section{Conclusion}

Computer vision combined with deep learning allows automated analysis of medical images.

However deployment in clinical environments requires:

\begin{itemize}
\item large validation datasets
\item interpretability
\item rigorous evaluation
\end{itemize}

\end{document}